% Research lemma supplied to the D.2.45 owner for independent audit.
% Not registered or asserted as a completed numbered theorem until audited.
\subsubsection*{An actual relative-characteristic computation of the refined class}

Let $X$ be a simply connected spin six-manifold with
$H_2(X;\mathbb Z)=H_3(X;\mathbb Z)=\mathbb Z/c$, where $c$ is even.
Put $U=X\setminus\{p\}$, with the point away from all the cycles
used below. Let $C$ be an oriented integral two-cycle, possibly a
sum of intersecting complex curves, and suppose that $TX$ has a
specified full reference frame $\rho$ on $U\setminus C$ whose
relative spin characteristic obstruction is
\[
                     \lambda(TX,\rho)=-\operatorname{PD}(C).
\tag{D.RP1}
\]
The equality is in the relative degree-four group for this
complement. Thus it specifies local obstruction coefficients,
not only the absolute characteristic class. The semistable
critical-cycle construction supplies precisely such a frame;
the original type-$II$ chart contributes its actual Jacobian
multiplicity two.

Let $N\subset U\setminus C$ be an oriented normally framed
three-manifold whose projection generates the two-primary
summand of $H_3(X;\mathbb Z)$, and put $d_N=\operatorname{ord}[N]$.
The class need not generate the odd-primary summand. Assume
that its stable tangent framing induced by $\rho$ and that
normal frame is an even multiple of $\nu$ in
$\Omega^{\rm fr}_3=\mathbb Z/24$. The ordinary regular three-torus
has this property by its Lie product framing and the explicitly
controlled determinant-character changes.

Choose an integral four-chain $B$ with $\partial B=d_NN$, made
transverse to the smooth strata of $C$ and disjoint from its
finitely many singular points. Define the actual rational link
\[
                         \ell(N,C)=\frac{B\cdot C}{d_N}.
\tag{D.RP2}
\]
It is a rational number, not merely a class modulo integers.
It is independent of $B$: the difference is a closed four-cycle,
and the homology class of $C$ is torsion. Here all of $H_2(X)$
is torsion. Changing the geometric representatives can change
the integer part, which is why their actual markings are retained.

We claim that Zhubr's refined class $P\in H^4(X;\mathbb Z)$ obeys
\[
                  \operatorname{lk}(N,\operatorname{PD}P)
                    =-\frac{\ell(N,C)}2\quad\text{in }\mathbb Q/\mathbb Z.
\tag{D.RP3}
\]
The linking pairing uses the orientation and spin obstruction
normalizations in (D.RP1). A simultaneous reversal of the
intersection convention changes both sides. In particular this
is not an instruction to divide an element of an even cyclic
group by two. Its right side first uses the specified rational
number from the actual cycles.

\paragraph{The admissible three-frame and a genuine embedded sphere.}
Write $d_N=2^r m$ with $m$ odd. Join $m$ parallel copies of $N$ by tubes with their
normal frames. The trace of the joining is a normally framed
embedded cobordism with one-handles. Apply Kaplan's prescribed
spin almost-framing construction to the resulting connected
three-manifold. Deleting its zero-handle and reversing gives a
spin cobordism to $S^3$ with only two-handles
\cite[Definition 1.7, Remark 1.8, and Theorem 3.1, pp.239, 243--246]{mapsKaplan79}.

These handles are realized in $U$, not just abstractly in a
seven-manifold. Their attaching circles bound disjoint embedded
disks with interior disjoint from the current three-manifold,
since $U$ is simply connected and $2+3<6$, $2+2<6$. A prescribed
boundary collar uses the first normal direction. The remaining
rank-four disk framing extends exactly when the specified spin
structure extends over the handle, since its obstruction is
$\pi_1SO(4)=\mathbb Z/2$. Thus the handle traces in $U\times I$
give a normally framed cobordism ending at an actual embedded
$S^3\subset U$ representing $m[N]$, a generator of the
two-primary summand.

Now choose the three-frame $v$ on $U$ supplied by Zhubr's
Theorem~1 for this particular outgoing embedded sphere
\cite[Theorem 1, p.6]{ffZhubr13}. The
theorem permits any such chosen embedded basis. Thus $v$ is
admissible on the actual sphere just constructed, not on an
unrelated representative. Put $\xi=v^\perp$. This oriented
rank-three bundle has a spin lift. Its primary obstruction to
a full frame is $o_4(\xi)=P$, and its stabilized spin
characteristic class is
\[
                              \lambda(\xi)=2o_4(\xi).
\tag{D.RP4}
\]
Choose a spin frame $\tau$ of $\xi|_N$, possible because
$B\operatorname{Spin}(3)$ is three-connected.

The spin frame $\tau$ of $\xi$ extends over this trace. Relative
to the incoming copies of $N$, the trace has only one- and
two-handles, and the corresponding extension groups for a spin
three-frame are $\pi_0\operatorname{Spin}(3)=0$ and
$\pi_1\operatorname{Spin}(3)=0$. Together with $v$ this gives a
full ambient tangent frame along the trace. The induced stable
tangent framings at its boundary are therefore bordant. On the
outgoing sphere, admissibility of $v$ means that the induced
stable framed class is even: the tangent Lie three-frame has
odd class $\nu$, the required relative three-frame is the
nonzero element of $\pi_3V_{6,3}=\mathbb Z/2$, and a change of
normal $SO(3)$ frame changes the stable class by an even multiple
of $\nu$. Hence
\[
                             [N]_{v+\tau}=0\pmod2,
\tag{D.RP5}
\]
because $m$ is odd. This argument uses neither a global full
frame on $U$ nor the vanishing of the odd-primary part of
$\lambda(TX)$.

\paragraph{Relative classes retain the missing integer.}
The complement $U\setminus C$ is simply connected: general
position moves any disk bounding a loop away from the real
codimension-four curve strata and their higher-codimension
singularities. Thus its tangent spin structure is unique, and
the frame $\rho$ induces the restriction of the spin structure
of $U$. Both $\rho|_N$ and $v+\tau$ are consequently compatible
spin frames. Their difference lifts to $\operatorname{Spin}(6)$.
Let $a\in H^3(N;\mathbb Z)$ be the pullback of its normalized
degree-three generator, and put $k=\langle a,[N]\rangle$.
On a three-dimensional complex a map to $\operatorname{Spin}(6)$
is determined by this class. The stable $J$ homomorphism sends
its generator to $\nu$, so (D.RP5) and the assumed even
reference framing imply that $k$ is even.

In $H^4(U,N;\mathbb Z)$ the two relative characteristic
classes satisfy, with one consistent choice of transition-map
direction,
\[
 \lambda(TX,\rho|_N)
       =2o_4(\xi,\tau)+\delta a.
\tag{D.RP6}
\]
This follows from the universal integral transgression of
$\lambda$ for $\operatorname{Spin}(6)$ and from (D.RP4).
It is an identity of relative classes with specified boundary
frames. Put $q=\langle o_4(\xi,\tau),B\rangle\in\mathbb Z$.
Pairing (D.RP6) with $B$ gives
\[
                        -B\cdot C=2q+d_Nk.
\tag{D.RP7}
\]
By the integral definition of the torsion linking pairing,
$q/d_N\bmod\mathbb Z$ is
$\operatorname{lk}(N,\operatorname{PD}P)$. Divide (D.RP7) as an
identity of integers, then use that $k$ is even. This proves
(D.RP3). Changing the spin frame $\tau$ changes $q$ by a multiple
of $d_N$ and $k$ by an even integer, as also follows from
$\pi_3\operatorname{Spin}(3)\to\pi_3\operatorname{Spin}(6)$
being multiplication by two.

For a cyclic group the pairing with $N$ determines the
two-primary component of $P$; it determines all of $P$ when
$[N]$ generates the entire group. The calculation may be reduced in $\mathbb Q/\mathbb Z$
after evaluating (D.RP2). For example, if $5C$ bounds, the
integral link with that literal cycle gives
$\ell(N,C)=\operatorname{lk}(N,5C)/5$; the odd denominator is
retained until this final reduction. No assertion that the full
critical cycle is nullhomologous is then required.
